\section{研究问题}

本研究关注全景室内场景中半自动角点标注的质量保障问题。在无法对所有样本进行全量精标的现实约束下，如何系统性地评估半自动标注的效率、质量与可靠性，并利用场景特异可靠度支持任务路由。

技术主线：采用可审计三段式规程 Pilot→PreScreen→Main；主实验前锁定阈值与规则，包括 MDE、$w_u$ 映射、关键超参数与扰动算子库；在 non-IID/OOD 条件下以标注前可得的代理信号（如 $d_t$ 与 $g_t$）驱动首轮风险分层与应激触发，并以标注后 meta-label 共识作诊断与审计；输出画像、矩阵与反例库等过程证据。
\subsection{研究问题 RQ}

本文围绕三个核心研究问题展开，分别对应效率、质量与路由三个维度。

\textbf{RQ1 效率评估}：半自动初始化是否降低真实活跃标注时间？下降比例是多少？
\\ 主终点：Manual vs Semi-Auto 的 \texttt{active\_time} 中位数差异（秒）及 95\% CI。 \texttt{active\_time} 的操作定义、采集实现与聚合口径见第~\ref{sec:method-active-time}节。
\\ 风险控制：分析 \texttt{active\_time} 与 $\mathrm{IoU}_{\mathrm{edit}}$ 的关系，排除时间缩短仅源于随意修改。

\textbf{RQ2 质量与可靠性}：在部分复核约束下，半自动流程是否改变标注一致性与可靠性分布，能否系统性暴露反例类型？
\\ 主终点：在预注册的同图配对子集（$N_{img}=25$）上比较 $\mathrm{IAA}_t$ 的变化，效应量取 $\mathrm{median}(\Delta_t)$，并使用配对置换检验与配对 bootstrap 给出 95\% CI；K-S 检验仅作为附录敏感性分析，不作为主终点。
\\ 反例类型：Type 1 低一致性；Type 2 异常编辑；Type 3 几何合法性失败（定义见第~\ref{sec:method-counterexample}节）。
\\ 可复现筛选规则：按第~\ref{sec:method-counterexample}节规则自动筛选候选，最终反例需专家复核与记录原因。

\textbf{RQ2b 加权共识}：加权共识作为干预手段，是否能在不引入循环论证的前提下提升共识稳定性？
\\ 评价口径：主评价使用独立参考子集或未加权一致性指标；加权共识的适用前提、锁定流程与报告方式见第~\ref{sec:method-weighted-consensus}节。

\textbf{RQ3 路由与优化}：在固定标注预算下，能否通过预筛选、工人画像与场景/分布感知路由，在 non-IID 与 OOD 场景中维持一致性并改善预算效率？同时，对门控失败与字段缺失等安全性事件保持可追溯披露，避免以静默过滤换取表面指标。
\\ 场景定义：以 \texttt{difficulty} 与 \texttt{model\_issue} 的共识标签作为场景代理。全部标签进入审计集合并在报告中给出频率、共现与反例分布。为避免小样本下过度细分导致数据稀疏，路由仅在不超过 $S_{\mathrm{core,max}}=4$ 个核心场景上启用，并必须披露启用率与退化率。
\\ 场景特异可靠度：在上述有限场景集合上估计标注者的场景内一致性中位数，并通过启用门槛、启用率/退化率披露与保守 LCB 排序保证稳健性；固定系数收缩仅作为附录敏感性分析报告（不作为主实现）。定义见第~\ref{sec:method-scene-reliability}节。
路由策略：Random；Global（仅用全局 $r_u$）；Full（PreScreen + 场景分层 + OOD 触发），定义与应激情景见第~\ref{sec:exp-iid-stress}节。
\\ 评估协议：离线模拟，在每任务已有的 k 份真实标注中，模拟策略分配任务并取实际提交标注作为输出，以避免因使用真实标签导致的信息泄露与乐观估计。
\\ 指标（主）：一致性代理提升、预算效率（序贯停止的 $\bar{k}_{\mathrm{used}}$ 或等价预算-覆盖关系）。
\\ 安全性/数据可用性审计（非主终点）：门控失败率与失败原因分布、关键字段缺失率；用于证明不静默过滤与口径 T/I/M 的可追溯性，不作为“路由收益”的主证据。


